%%%
%%% Radical "⽸"
%%%
\section*{Radical 121: ``⽸''}\addcontentsline{toc}{section}{Radical 121: ⽸}\addcontentsline{loh}{figure}{\#\#\#\# 121: ⽸}

%%%%%%%%%% 缸 %%%%%%%%%%
\subsection*{缸}\addcontentsline{loh}{figure}{缸}

\begin{Entry}{缸}{9}{⽸}
  \begin{Phonetics}{缸}{gang1}[][HSK 7-9]
    \definition[口,个]{s.}{jarra; pote de barro; recipientes feitos de barro, porcelana, vidro, etc. geralmente têm uma boca grande e um fundo pequeno | composto de areia, argila, etc. para fazer louça de barro | vaso em forma de jarro; objetos em forma de potes}
  \end{Phonetics}
\end{Entry}

%%%%%%%%%% 缺 %%%%%%%%%%
\subsection*{缺}\addcontentsline{loh}{figure}{缺}

\begin{Entry}{缺}{10}{⽸}
  \begin{Phonetics}{缺}{que1}[][HSK 3]
    \definition{adj.}{incompleto; imperfeito}
    \definition[种]{s.}{vaga; abertura; falta}
    \definition{v.}{estar com falta de; faltar | estar ausente}
  \synonymref{残}{can2}
  \synonymref{短}{duan3}
  \synonymref{亏}{kui1}
  \antonymref{满}{man3}
  \antonymref{欠}{qian4}
  \antonymref{全}{quan2}
  \antonymref{余}{yu2}
  \antonymref{圆}{yuan2}
  \end{Phonetics}
\end{Entry}

\begin{Entry}{缺口}{10,3}{⽸,⼝}
  \begin{Phonetics}{缺口}{que1kou3}[][HSK 7-9]
    \definition[个]{s.}{fenda; entalhe; brecha; reentrância; uma lacuna formada pela falta de uma parte de um objeto | déficit; falta de fundos, materiais, etc.; suprimentos, fundos, etc. insuficientes}
  \seealsoref{缺口儿}{que1kou3r5}
  \end{Phonetics}
\end{Entry}

\begin{Entry}{缺口儿}{10,3,2}{⽸,⼝,⼉}
  \begin{Phonetics}{缺口儿}{que1kou3r5}
    \definition{s.}{brecha; lacuna}
  \end{Phonetics}
\end{Entry}

\begin{Entry}{缺乏}{10,4}{⽸,⼃}
  \begin{Phonetics}{缺乏}{que1fa2}[][HSK 5]
    \definition{v.}{faltar; estar em falta de; não ter ou não ter totalmente (algo que deveria possuir ou é desejaria possuir)}
  \synonymref{不足}{bu4zu2}
  \synonymref{短缺}{duan3que1}
  \synonymref{欠缺}{qian4que1}
  \synonymref{缺少}{que1shao3}
  \synonymref{缺失}{que1shi1}
  \antonymref{充分}{chong1fen4}
  \antonymref{充足}{chong1zu2}
  \antonymref{丰富}{feng1fu4}
  \antonymref{富有}{fu4you3}
  \antonymref{足够}{zu2gou4}
  \end{Phonetics}
\end{Entry}

\begin{Entry}{缺少}{10,4}{⽸,⼩}
  \begin{Phonetics}{缺少}{que1shao3}[][HSK 3]
    \definition{v.}{falta; estar com falta de; estar em falta de; geralmente se refere à falta de pessoas ou coisas}
  \seealsoref{缺乏}{que1fa2}
  \synonymref{短缺}{duan3que1}
  \synonymref{短少}{duan3shao3}
  \synonymref{欠缺}{qian4que1}
  \antonymref{充分}{chong1fen4}
  \antonymref{充足}{chong1zu2}
  \antonymref{多余}{duo1yu2}
  \antonymref{具备}{ju4bei4}
  \antonymref{齐备}{qi2bei4}
  \antonymref{剩余}{sheng4yu2}
  \antonymref{拥有}{yong1you3}
  \end{Phonetics}
\end{Entry}

\begin{Entry}{缺失}{10,5}{⽸,⼤}
  \begin{Phonetics}{缺失}{que1shi1}[][HSK 7-9]
    \definition{s.}{defeito; desvantagem; deficiência}
    \definition{v.}{perder; faltar; ter pouca}
  \end{Phonetics}
\end{Entry}

\begin{Entry}{缺点}{10,9}{⽸,⽕}
  \begin{Phonetics}{缺点}{que1dian3}[][HSK 3]
    \definition[个,些]{s.}{desvantagem; deficiência; inconveniência; ponto fraco; uma deficiência ou imperfeição}
  \seealsoref{缺陷}{que1xian4}
  \synonymref{弊端}{bi4duan1}
  \synonymref{错误}{cuo4wu4}
  \synonymref{短处}{duan3chu5}
  \synonymref{过失}{guo4shi1}
  \synonymref{坏处}{huai4chu5}
  \synonymref{漏洞}{lou4dong4}
  \synonymref{毛病}{mao2bing4}
  \synonymref{弱点}{ruo4dian3}
  \antonymref{长处}{chang2chu4}
  \antonymref{亮点}{liang4dian3}
  \antonymref{优点}{you1dian3}
  \end{Phonetics}
\end{Entry}

\begin{Entry}{缺席}{10,10}{⽸,⼱}
  \begin{Phonetics}{缺席}{que1//xi2}[][HSK 7-9]
    \definition{v.+compl.}{ausentar"-se; estar ausente (de uma reunião, etc.)}
  \end{Phonetics}
\end{Entry}

\begin{Entry}{缺陷}{10,10}{⽸,⾩}
  \begin{Phonetics}{缺陷}{que1xian4}[][HSK 6]
    \definition[个,处,项]{pron.}{defeito; falha; inconveniência; mancha; um lugar onde uma pessoa ou coisa está incompleta ou tem falhas porque algo está faltando}
  \seealsoref{缺点}{que1dian3}
  \synonymref{弊端}{bi4duan1}
  \synonymref{短处}{duan3chu5}
  \synonymref{毛病}{mao2bing4}
  \synonymref{弱点}{ruo4dian3}
  \antonymref{长处}{chang2chu4}
  \antonymref{健全}{jian4quan2}
  \antonymref{亮点}{liang4dian3}
  \antonymref{特长}{te4chang2}
  \antonymref{完美}{wan2mei3}
  \antonymref{完善}{wan2shan4}
  \antonymref{优点}{you1dian3}
  \antonymref{优势}{you1shi4}
  \end{Phonetics}
\end{Entry}

\begin{Entry}{缺勤}{10,13}{⽸,⼒}
  \begin{Phonetics}{缺勤}{que1//qin2}
    \definition{v.+compl.}{ausentar"-se do serviço; ausentar"-se do trabalho; faltar à escola}
  \seealsoref{旷工}{kuang4//gong1}
  \synonymref{缺席}{que1//xi2}
  \end{Phonetics}
\end{Entry}

%%%%%%%%%% 罐 %%%%%%%%%%
\subsection*{罐}\addcontentsline{loh}{figure}{罐}

\begin{Entry}{罐}{23}{⽸}
  \begin{Phonetics}{罐}{guan4}[][HSK 7-9]
    \definition{clas.}{lata; jarra; gavetas e recipientes de água feitos de cerâmica ou metal}[我买了一罐可乐。===Comprei uma lata de Coca"-Cola.]
    \definition{s.}{lata; jarra; jarro; pote; tanque | cuba de carvão; vagão de caçamba para carregamento de carvão em minas de carvão}
  \end{Phonetics}
\end{Entry}

\begin{Entry}{罐头}{23,5}{⽸,⼤}
  \begin{Phonetics}{罐头}{guan4tou5}[][HSK 7-9]
    \definition[个,盒,瓶]{s.}{lata; jarra | enlatado; comida enlatada é a abreviação de 罐头食品, que é processada e embalada em latas de ferro seladas ou garrafas de vidro e pode ser armazenada por um longo tempo}
  \seealsoref{罐头食品}{guan4tou2 shi2pin3}
  \end{Phonetics}
\end{Entry}

\begin{Entry}{罐头食品}{23,5,9,9}{⽸,⼤,⾷,⼝}
  \begin{Phonetics}{罐头食品}{guan4tou2 shi2pin3}
    \definition{s.}{alimentos enlatados; produtos enlatados}
  \end{Phonetics}
\end{Entry}

%%%%% EOF %%%%%

